\chapter{Phase 6: Multi-player mode (future)}

\section{Vision}

The configuration described up to the previous chapter supports a single phone. The original television experience, however, contemplated multiple players calling at the same time, listening to the game's audio while waiting, and receiving rotating 30-to-60-second turns.

This phase describes the design needed to scale the system to that full experience. \textbf{It is not documented as a step-by-step executable guide} (that is work for a 2.0 version of the project), but as an architectural roadmap.

\section{Hardware changes}

\begin{itemize}
  \item Replace the \textbf{HT801} (1 FXS port) with:
    \begin{itemize}
      \item \textbf{HT802}: 2 FXS ports. Approx. cost \$55,000 CLP.
      \item \textbf{HT814}: 4 FXS ports. Approx. cost \$110,000 CLP.
      \item \textbf{Cisco SPA8000} (used): 8 FXS ports. Approx. cost \$80,000--120,000 CLP.
    \end{itemize}
  \item Acquire as many additional analog phones as FXS ports you'll use.
  \item Optionally, a USB HDMI capture card to stream or record matches.
\end{itemize}

\section{Asterisk changes}

\subsection{pjsip.conf: new extensions}

Add extensions 102, 103, 104 according to the number of FXS ports, replicating the \texttt{endpoint}/\texttt{auth}/\texttt{aor} block that 101 uses (inheriting from \texttt{endpoint-template}).

\subsection{confbridge.conf: conference rooms}

Create the file \texttt{/etc/asterisk/confbridge.conf} with profiles for two kinds of participants:

\begin{astcode}
[default_bridge]
type=bridge
max_members=10
mixing_interval=20

[player_active]
type=user
admin=no
marked=yes
quiet=yes

[player_waiting]
type=user
admin=no
marked=no
quiet=yes
music_on_hold_when_empty=no
\end{astcode}

\subsection{extensions.conf: queue logic}

Modify extension \texttt{9000} so that the first participant enters as \emph{player\_active} and subsequent ones as \emph{player\_waiting}:

\begin{astcode}
[hugo-game]
exten => 9000,1,Answer()
 same => n,Set(POSITION=${CONFBRIDGE_INFO(parties,hugo-room)})
 same => n,GotoIf($[${POSITION} = 0]?active:waiting)
 same => n(active),ConfBridge(hugo-room,default_bridge,player_active)
 same => n(waiting),Playback(esperando-turno)
 same => n,ConfBridge(hugo-room,default_bridge,player_waiting)
\end{astcode}

\section{Turn logic}

An additional component is needed (it could be another Rust process or a module inside the existing bridge) that every 60 seconds:

\begin{enumerate}
  \item Identifies the current \emph{player\_active} via the \texttt{ConfbridgeList} event or \texttt{ConfbridgeList} queries over AMI.
  \item Demotes them to \emph{waiting} via \texttt{ConfbridgeMute} and a profile reconfiguration.
  \item Promotes the next player in queue (FIFO) to \emph{active}.
  \item Announces the turn change by playing a pre-recorded message to everyone.
\end{enumerate}

Only the active player should be able to control the character. This is achieved by having the bridge filter \texttt{UserEvent} messages by the \texttt{Channel} field and only accept those from the active channel.

\section{Game audio in the conference}

For waiting players to hear Hugo, the audio from DOSBox must be routed into the conference room.

\subsection{Path A: ALSA loopback}

\begin{shellcode}
sudo modprobe snd-aloop
\end{shellcode}

This creates a virtual device \texttt{hw:Loopback,0}. DOSBox is configured to output audio to that device. Asterisk opens the other end of the loopback as an input channel and \emph{bridges} it into the room.

\subsection{Path B: PipeWire}

More flexible and modern but with higher resource usage. On the Pi 4 (2\,GB) it could be a bottleneck; on Pi 5 it works without trouble.

\section{Physical topology}

\begin{center}
\begin{tikzpicture}[
  every node/.style={font=\small},
  box/.style={draw=hugoblue, thick, rounded corners=3pt, fill=hugoblue!10, align=center}
]
  \node[box, minimum width=2cm, minimum height=1cm] (tv) at (0,2.5) {Television};
  \node[box, minimum width=3cm, minimum height=1.2cm] (pi) at (0,0) {Raspberry Pi 4};
  \node[box, minimum width=3cm, minimum height=1.2cm] (ata) at (0,-2.5) {HT814 (4 FXS)};

  \draw[->, thick, hugoaccent] (pi) -- node[right,font=\tiny] {HDMI} (tv);
  \draw[<->, thick, hugoaccent] (pi) -- node[right,font=\tiny] {Ethernet} (ata);

  \node[box, minimum width=1.5cm, minimum height=0.8cm] (p1) at (-4.5,-2.5) {Phone 1};
  \node[box, minimum width=1.5cm, minimum height=0.8cm] (p2) at (-2.5,-4.5) {Phone 2};
  \node[box, minimum width=1.5cm, minimum height=0.8cm] (p3) at (2.5,-4.5) {Phone 3};
  \node[box, minimum width=1.5cm, minimum height=0.8cm] (p4) at (4.5,-2.5) {Phone 4};

  \draw[-, hugoaccent] (p1) -- (ata);
  \draw[-, hugoaccent] (p2) -- (ata);
  \draw[-, hugoaccent] (p3) -- (ata);
  \draw[-, hugoaccent] (p4) -- (ata);
\end{tikzpicture}
\end{center}

\section{Optional features}

\subsection{Player identification}

At the start of every call, before entering the room, the player can be asked to dial a digit to identify themselves:

\begin{plaincode}
"Welcome to Hugo. Press 1 if you are new, or your player number
if you already have one."
\end{plaincode}

This makes it possible to keep a persistent ranking across sessions.

\subsection{Scoreboard / ranking}

A local SQLite database on the Pi can store scores per player. Each time Hugo registers a life lost or a level completed (information that could be captured via OCR on the HDMI output or by reading the DOSBox process memory), the scoreboard is updated.

\subsection{On-screen scoreboard display}

An overlay can be drawn on top of the HDMI output via:

\begin{itemize}
  \item OBS Studio in composition mode.
  \item A second process drawing with SDL2 on a transparent layer.
  \item A small secondary LCD screen dedicated to the scoreboard.
\end{itemize}

\subsection{Streaming to Twitch or YouTube}

With a USB HDMI capture card plugged into the Mac (not the Pi, to avoid overloading it), matches can be broadcast. This lets other friends watch the game without being physically present.

\section{Effort estimate}

\begin{table}[h]
\centering
\begin{tabular}{lc}
\toprule
\textbf{Component} & \textbf{Approx. effort} \\
\midrule
ConfBridge and new SIP extensions & 4 hours \\
Turn logic (Rust) & 16 hours \\
ALSA audio loopback & 4 hours \\
Player identification & 4 hours \\
SQLite persistence & 4 hours \\
Scoreboard overlay & 12 hours \\
Tests and calibration & 8 hours \\
\midrule
\textbf{Estimated total} & \textbf{$\sim$52 hours} \\
\bottomrule
\end{tabular}
\caption{Time estimate for implementing multi-player mode.}
\end{table}

\section{Wrap-up}

The base system designed in chapters 4--8 is intended to grow into this multi-player version without rewriting components. When the time comes to implement it, a new detailed design phase should be opened taking these foundations as a reference.
